% ============================================================
% Chapter 1: Introduction
%   Section: Scope
%     Subsection: In Scope
% Included from chapters/01-introduction/06-scope/scope.tex
% ============================================================

\subsection{In Scope}

\begin{itemize}
    \item \textbf{Discovery and selection} of articles from 33 enabled
    news sources (20 in English, 13 in Spanish) grouped into five topic
    groups --- society, science, environment, culture and health ---
    with a screen for the publication's mission, a ranking of the
    candidates and a choice by a person or a language model.
    \item \textbf{Extraction} of article text and metadata through a
    cascade of strategies, for the articles analysed and for every web
    page read as evidence.
    \item \textbf{Enrichment} with local models: entities, keywords,
    topics, sentiment, editorial quality, candidate claims and
    embeddings.
    \item \textbf{Fact-checking}: admission, claim selection, open-web
    evidence retrieval, ranking with a pertinence gate, verification by
    a local language model, confidence recalibration and article-level
    aggregation.
    \item \textbf{Storage}: a three-layer data lake with lineage, a
    vector store and a graph database written by every analysis.
    \item \textbf{Evaluation}: a pilot on X-Fact, a hand-labelled
    validation set from the system's domain, an evaluation harness, and
    experiments on the stages before the verdict.
    \item \textbf{A reader view} of the published articles, and a
    deployment to Google Cloud designed and checked against a local
    emulator.
\end{itemize}
